\clearpage
\VerdierPhysicalPageAnchor{0176}%
\begingroup\small\setlength{\parskip}{0pt}\setlength{\emergencystretch}{1.5em}
\setlength{\abovedisplayskip}{4pt plus 2pt minus 2pt}%
\setlength{\belowdisplayskip}{4pt plus 2pt minus 2pt}%

\VerdierSemanticAnchor{ch3d-1-2}
\subsection*{1.2. 导出范畴的若干重要子范畴。}

\VerdierSemanticAnchor{ch3d-1-2-1}
\noindent\textbf{命题 1.2.1.}\quad
\textit{a) 子范畴 $K^{,*}(\mathcal A)\subset K(\mathcal A)$ 是右、左
$\operatorname{Ac}(\mathcal A)$-局部化的（第二章，2.3.5）。}

\noindent\textit{b) 子范畴 $K^{+,*}(\mathcal A)\subset K(\mathcal A)$
是右 $\operatorname{Ac}(\mathcal A)$-局部化的。}

\noindent\textit{b)\textsuperscript{bis} 子范畴
$K^{-,*}(\mathcal A)\subset K(\mathcal A)$ 是左
$\operatorname{Ac}(\mathcal A)$-局部化的。}

\noindent\textit{c) 子范畴 $K^{b}(\mathcal A)\subset K^{+}(\mathcal A)$
是左 $\operatorname{Ac}^{+}(\mathcal A)$-局部化的。}

\noindent\textit{c)\textsuperscript{bis} 子范畴
$K^{b}(\mathcal A)\subset K^{-}(\mathcal A)$ 是右
$\operatorname{Ac}^{-}(\mathcal A)$-局部化的。}

\noindent\textit{d) $K^{,+}(\mathcal A)$ 的每个对象都有一个到
$K^{+}(\mathcal A)$ 某个对象的拟同构。}

\noindent\textit{d)\textsuperscript{bis} $K^{,-}(\mathcal A)$ 的每个对象
都从 $K^{-}(\mathcal A)$ 的某个对象接收一个拟同构。}

\noindent\textit{e) $K^{,b}(\mathcal A)$ 的每个对象都有一个到
$K^{+,b}(\mathcal A)$ 某个对象的拟同构。}

\noindent\textit{e)\textsuperscript{bis} $K^{,b}(\mathcal A)$ 的每个对象
都从 $K^{-,b}(\mathcal A)$ 的某个对象接收一个拟同构。}

\noindent\textit{f) $K^{+,b}(\mathcal A)$ 的每个对象都从
$K^{b}(\mathcal A)$ 的某个对象接收一个拟同构。}

\noindent\textit{f)\textsuperscript{bis} $K^{-,b}(\mathcal A)$ 的每个对象
都有一个到 $K^{b}(\mathcal A)$ 某个对象的拟同构。}

\VerdierSemanticAnchor{ch3d-1-2-1-truncations}
对每个复形 $X$，以 $X(n,+\infty)$（相应地，$X(-\infty,n)$）表示复形
\[
X(n,+\infty)=\cdots\longrightarrow0\longrightarrow0
\longrightarrow X^n/\operatorname{Im}(d_X^{n-1})
\longrightarrow X^{n+1}\longrightarrow\cdots
\]
（相应地，
\[
X(-\infty,n)=\cdots\longrightarrow X^{n-3}\longrightarrow X^{n-2}
\longrightarrow\operatorname{Ker}(d_X^{n-1})
\longrightarrow0\longrightarrow0\longrightarrow\cdots
\]）
并以 $p\colon X\longrightarrow X(n,+\infty)$（相应地，
$i\colon X(-\infty,n)\longrightarrow X$）表示如下复形态射：
\[
\begin{aligned}
p^q&=\operatorname{id},&&q>n,\\
p^n&\colon X^n\longrightarrow X^n/\operatorname{Im}(d_X^{n-1})
  &&\text{（典范态射）},\\
p^q&=0,&&q\leq n-1.
\end{aligned}
\]
\endgroup
